% =====================================================================
%  Chapitre 4 - Sprint 1 : identites, acces et securite
% =====================================================================
\chapter{Sprint 1 : identites, acces et securite}
\label{chap:sprint1}

\section*{Introduction}
\addcontentsline{toc}{section}{Introduction}

Ce premier sprint construit le socle sur lequel reposent tous les
suivants : sans identites ni habilitations, aucun ecran de la plateforme
ne peut etre protege et aucune action ne peut etre imputee a son auteur.
Il couvre l'inscription, l'authentification, la gestion du compte
personnel et l'administration des comptes par les responsables.

% =====================================================================
\section{Planification du sprint}

\subsection{Objectifs}

\begin{enumerate}
  \item Permettre a un visiteur de creer un compte et d'en confirmer
    l'adresse electronique.
  \item Authentifier les utilisateurs et proteger les acces par un second
    facteur optionnel.
  \item Offrir au chauffeur un mode de connexion adapte a la mobilite.
  \item Donner a l'administrateur les moyens de statuer sur les demandes
    d'inscription et de gerer le cycle de vie des comptes.
  \item Rendre toute connexion tracable.
\end{enumerate}

\subsection{Carnet du sprint}

\begin{table}[H]
\centering
\caption{Carnet du sprint 1}
\label{tab:backlog-s1}
\small
\begin{tabularx}{\textwidth}{@{}cp{2.5cm}Xcc@{}}
\toprule
\textbf{N\textsuperscript{o}} & \textbf{Acteur} & \textbf{Besoin exprime} &
\textbf{Prio.} & \textbf{Pts} \\ \midrule
1 & Visiteur & Creer un compte & Must & 5 \\
2 & Visiteur & Confirmer son adresse electronique & Must & 3 \\
3 & Utilisateur & Se connecter a son espace & Must & 5 \\
4 & Utilisateur & Activer un second facteur & Should & 8 \\
5 & Chauffeur & Se connecter par code PIN & Should & 5 \\
6 & Utilisateur & Gerer son profil et son mot de passe & Must & 3 \\
7 & Administrateur & Statuer sur les demandes d'inscription & Must & 5 \\
8 & Responsable & Gerer les habilitations d'un compte & Must & 3 \\
9 & Administrateur & Archiver un compte sans perdre son historique & Must & 5 \\
10 & Administrateur & Consulter l'historique des connexions & Should & 3 \\
\midrule
\multicolumn{4}{@{}r}{\textbf{Total}} & \textbf{45} \\
\bottomrule
\end{tabularx}
\end{table}

% =====================================================================
\section{Elaboration technique}

\subsection{Modelisation fonctionnelle}

\figReelle{iam-usecase.png}
  {Cas d'utilisation du service des identites et des acces}
  {fig:uc-iam}{0.50}

Le diagramme de la figure~\ref{fig:uc-iam} fait apparaitre trois
regroupements. \essentiel{L'acces au service} rassemble ce qui precede
l'entree dans l'application, \essentiel{le compte personnel} ce que
l'utilisateur gere pour lui-meme, et \essentiel{l'administration des
comptes} ce qui est reserve au responsable et a l'administrateur.

Deux relations d'extension meritent d'etre relevees. \exemple{Valider un
second facteur} etend \exemple{se connecter} : l'etape n'a lieu que si
l'utilisateur a active la protection. \exemple{Archiver un compte} etend
\exemple{supprimer un compte} : l'archivage est propose lorsque le compte
porte un historique metier. Le cas \exemple{consulter les dependances
metier} est en revanche \essentiel{inclus}, la verification etant
systematique avant toute suppression.

\begin{table}[H]
\centering
\caption{Description textuelle du cas « se connecter »}
\label{tab:desc-connexion}
\small
\begin{tabularx}{\textwidth}{@{}p{2.5cm}X@{}}
\toprule
\textbf{Rubrique} & \textbf{Contenu} \\ \midrule
Acteur & Utilisateur \\
Precondition & Le compte existe et son statut est ACTIF \\
Scenario nominal & Saisie des identifiants, verification de l'empreinte,
  emission d'un jeton d'acces valable soixante minutes et d'un jeton de
  rafraichissement, journalisation de la connexion \\
Extension & Si le second facteur est actif, un jeton temporaire est emis
  et le code est demande avant la delivrance du jeton d'acces \\
Alternatives & Identifiants invalides : compteur d'echecs incremente,
  verrouillage de quinze minutes au cinquieme echec. Compte en attente,
  suspendu ou archive : refus motive \\
\bottomrule
\end{tabularx}
\end{table}

\subsection{Conception statique}

\figReelle{iam-class-domain.png}
  {Modele du domaine du service des identites et des acces}
  {fig:dom-iam}{0.52}

Le modele repose sur une \essentiel{hierarchie d'heritage}. La classe
\exemple{Utilisateur} porte ce qui est commun ; \exemple{Client} et
\exemple{Employe} la specialisent, et \exemple{Chauffeur} specialise a son
tour \exemple{Employe}. Cette hierarchie est stockee dans une table
unique, distinguee par un discriminant : les specialisations partagent le
meme identifiant, ce qui evite une jointure a chaque lecture.

Le \essentiel{statut} du compte est un attribut a part entiere, distinct
de l'indicateur d'activite. Il decrit un cycle de vie en cinq etats : en
attente, actif, suspendu, desactive, archive. \essentiel{L'archivage n'est
pas une suppression} : la ligne est conservee, ce qui preserve les
missions, les factures et les interventions qui la referencent, et
l'operation reste reversible.

\subsection{Conception dynamique}

\figReelle{iam-seq-inscription.png}
  {Sequence : inscription d'un compte}
  {fig:seq-inscription}{0.72}

La sequence de la figure~\ref{fig:seq-inscription} montre que le compte
est cree au statut \essentiel{en attente} : il reste inutilisable tant
qu'un administrateur ne l'a pas approuve. Deux effets accompagnent la
creation : la publication d'un evenement sur le bus et l'envoi d'un
courriel de confirmation valable vingt-quatre heures.

\figReelle{iam-seq-approbation-compte.png}
  {Sequence : approbation d'un compte par l'administrateur}
  {fig:seq-approbation}{0.72}

L'approbation, figure~\ref{fig:seq-approbation}, fait passer le compte a
l'etat actif et publie a son tour un evenement. La lecture des deux
diagrammes met en evidence une constante de conception :
\essentiel{le controleur ne s'adresse jamais directement au depot}. Il
passe par une interface de service, satisfaite par une implantation qui
seule connait la base. C'est la regle de dependance exposee au
chapitre~\ref{chap:architecture}.

\subsection{Integration avec les autres services}

\figReelle{iam-communication.png}
  {Echanges du service des identites avec le reste de la plateforme}
  {fig:com-iam}{0.62}

Le service des identites est le seul a \essentiel{publier} des evenements
types sur le bus de messages. Le service de messagerie les consomme afin
de creer automatiquement les conversations d'accueil. Ce couplage est
\essentiel{asynchrone} : si la messagerie est indisponible, l'inscription
aboutit malgre tout et l'evenement sera traite au redemarrage.

% =====================================================================
\section{Realisation et tests de validation}

\figTrois{s1-connexion.jpg}{s1-double-facteur.jpg}{s1-connexion-pin.jpg}
  {Parcours d'authentification : connexion, validation du second facteur,
   connexion par code PIN du chauffeur}
  {fig:ecrans-auth}{0.40}

\figDeux{s1-comptes-attente.jpg}{s1-comptes-archives.jpg}
  {Administration des comptes : demandes en attente et comptes archives
   restaurables}
  {fig:ecrans-admin}{0.55}

\begin{table}[H]
\centering
\caption{Tests de validation du sprint 1}
\label{tab:tests-s1}
\small
\begin{tabularx}{\textwidth}{@{}cXp{3.0cm}c@{}}
\toprule
\textbf{N\textsuperscript{o}} & \textbf{Cas teste} & \textbf{Attendu} &
\textbf{Etat} \\ \midrule
T1 & Identifiants valides & Jeton delivre, connexion journalisee & Reussi \\
T2 & Cinq echecs consecutifs & Verrouillage de quinze minutes & Reussi \\
T3 & Connexion sur un compte en attente & Refus motive & Reussi \\
T4 & Second facteur actif & Code demande avant le jeton & Reussi \\
T5 & Archivage d'un compte porteur de missions & Missions preservees & Reussi \\
T6 & Restauration d'un compte archive & Compte de nouveau actif & Reussi \\
\bottomrule
\end{tabularx}
\end{table}

% =====================================================================
\section*{Conclusion}
\addcontentsline{toc}{section}{Conclusion}

Ce sprint livre le socle d'identites de la plateforme : inscription
soumise a approbation, authentification protegee par un second facteur
optionnel, connexion par code PIN pour les chauffeurs, gestion des
habilitations et archivage reversible. Les quarante-cinq points du carnet
ont ete realises et les six tests de validation sont passes.

Le choix le plus structurant de ce sprint n'est pas technique mais
metier : \essentiel{ne jamais supprimer un compte qui porte un
historique}. Il conditionne tous les modules suivants, ou les missions,
les factures et les interventions referencent leur auteur.

Le sprint suivant construit le second socle, celui des ressources
materielles : le referentiel de la flotte et des equipements.
